\chapter{为什么一个最大公因数能阻止所有平均}\label{sec:obstruction}
\section{七个数模七都等于一，会发生什么}
考虑三平均的配置 $(1^6,-6)$。总和是零，但每个数模 $7$ 都等于 $1$。任取三个，其和模 $7$ 等于 $3$；除以 $3$ 后，余数仍是 $1$。所以所有位置在每一步都模 $7$ 等于 $1$，不可能全变成零。

可能你会担心：平均后出现分数，还能谈模 $7$ 吗？可以，因为这些分母只是 $3$ 的幂，模 $7$ 都有逆元。例如 $1/3\equiv5\pmod7$，因为 $3\cdot5\equiv1$。所以“除以三”在模七下仍然说得清。

这个解释也说明为什么模平均元数 $p$ 不一样：$p$ 没有模 $p$ 的逆元。因此同样的阻碍不能拿来禁止 $G$ 中的 $p$ 因子。

\section{把例子写成一般证明}
先确认上一讲的归一化。
\InsertProof{lem:affine}
如果 $\ell\mid G$，意思正是所有坐标模 $\ell$ 相等。为什么共同余数不能恰好是零？因为那会使所有整数坐标都有公因子 $\ell$，违背本原性。这一句解释了归一化为何必须约掉共同因子。
\InsertProof{lem:obstruction}
其中 $G\mid n$ 很有用：将来要防范的素数不是所有素数，只是 $n$ 的非 $p$ 素因子。

称当前整数零和配置\textbf{合法}，是指它在每个 $\ell\mid n,\ell\ne p$ 处都不全同余。称保持这项性质的整数平均为\textbf{安全操作}。记这些需要防范的不同素数总数为
\[d_p(n)=\#\{\ell:\ell\mid n,\ell\ne p\text{ 为素数}\}.
\]
初始本原配置合法，当且仅当 $G$ 是 $p$ 的幂。操作后的配置未必还本原，所以后面会区分“当前实际整数值的余数”与“约去公因子后的判据”。

\section{合法输入也会被错误的一步带坏}
三平均七元输入 $(3^5,1,-16)$ 的总和为零，$G=\gcd(2,19)=1$。它是合法的。但如果先平均 $(3,1,-16)$，均值为 $-4$，输出就成了
\[(3^4,(-4)^3).
\]
现在所有数模 $7$ 都等于 $3$，再也不能归零。这说明 $G$ 并不是每一步自动保持的量。充分性不能任意平均，必须安排操作以保留必要的余数差异。

\section{为什么恰好少一个位置就会失败}
要证明 $2p+1$ 这个界不能再低，只需在 $2p$ 个位置找一个 $G=1$ 的失败例子。候选是
\[(-p,0^{p-1},1^p).
\]
它含 $0,1$，所以 $G=1$。如果只允许整数平均，唯一非恒等选择是 $(-p,0^{p-1})$，结果为 $((-1)^p,1^p)$。由于 $p$ 奇，这个配置所有数模 $2$ 都等于 $1$，失败。

但这里有个必须补上的问题：为什么不能先做一次产生分数的操作，绕过去？下面用“分母里有几个 $p$”排除它。

对非零有理数写 $x=p^h u/v$，$p\nmid uv$，定义 $v_p(x)=h$。例如 $v_3(18)=2$，$v_3(2/9)=-2$，并规定 $v_p(0)=+\infty$。只需要知道：两项赋值不同时，和的赋值等于较小的那个，因为提取最小 $p$ 幂后，一个单位加一个被 $p$ 整除的数不会变成零模 $p$。

首次产生分数时，恰好有新选中的 $p$ 个位置具有最低负赋值。下一次若只选它们中的 $k$ 个，$1\le k<p$，因 $p\nmid k$，最低项无法被其他较高赋值的项抵消；除以 $p$ 后，负赋值反而再降低一。这就是一直跟踪的东西。
\InsertProof{lem:fraction}
因此整数分支和分数分支都已排除，反例才真正完整。
\InsertProof{prop:lower}

\section{二平均为什么单列}
二平均的最终阈值是 $4$，而不是把 $p=2$ 代入奇素数公式所得的 $5$。文献\cite{mixing}的条件MC说：若所有整数输入在一个奇模数下相等，则目标均值也必须在那个余数类里。对本原零和输入，这正好等价于 $G$ 没有奇素因子。

三元反例为 $(-3,1,2)$。任意首次平均后都成 $(a,a,-2a)$、$a\ne0$；以后非恒等平均仍是同形状，只把 $a$ 变为 $-a/2$。它越变越小但不会是零。四元以上的充分性引用上述论文；其两个输入、两个输出的混合器按时间顺序执行，就是本题两个原位置同时取平均。
